
We began by asking whether better AI makes users intellectually lazier. For returning WildChat-1M users in 2023--2024, the empirical answer is no: prompt complexity grows significantly ($\tau$ = +0.744, p = 0.0005), not declines. But the appropriate comparison group has not been assembled, the measurement infrastructure to build that comparison is not publicly available, and the causal chain from AI quality improvement to behavioral change has not been verified. The result refutes the BAA directional prediction for prompt length in this cohort; it does not close the BAA empirical question.

Four contributions follow from this work: the first large-scale empirical test of the BAA disengagement directional prediction in public AI interaction logs (to our knowledge, pending citation verification; negative result); a validated behavioral trend measurement pipeline (Hamed-Rao Mann-Kendall on returning-user cohorts from WildChat-1M); documentation of a binding LMSYS data access constraint; and characterization of the selection bias problem with a proposed remedy.

The most critical next steps are: (1) a two-group comparison of returning versus non-returning users in a difference-in-differences design; (2) LMSYS primary research access for vote entropy analysis; (3) inclusion of a per-interaction AI quality covariate (e.g., model version or quality rating); and (4) implicit correction proxy development to replace the insufficient regex-based approach. The question of whether AI improvement changes human behavior --- and in what direction --- is answerable from existing interaction data. This study establishes the methodological foundation and documents the empirical constraints; resolving the ambiguity requires better data access, a more careful cohort design, and per-interaction AI quality measures not currently available in public datasets.

